% Proposed conceptual-article insertion only; existing manuscript untouched.
% Inherit Noto Sans/Fira Math from the selected master.
\section{A larger model still owes us a learning experiment}

A model could predict a hand's path more accurately while teaching its owner nothing. That distinction matters as the programme expands from recording-specific mappings toward a proposed Large Movement Model. The new hypothesis is that greater amounts and greater diversity of qualified neural, muscle and movement data could improve generalization. Brain2Body asks what any resulting prediction would let a trainee do, what information or assistance supplied the gain, and what remains when teaching ends. A scaling result would answer the first modelling question only within its tested domain.

\subsection{Four training sources, five component roles}

The expanded training proposal has at least four source groups: neural observations; muscle and other myographic observations; measured motion and displacement; and data generated by physics-based motion or musculoskeletal simulation. Suitable neural sources may span brain, spinal, peripheral and relevant visceral recordings. Their actual modalities, sites and task roles remain distinct. The present emphasis is motor--sensory. Myoception is the programme's term for the sensory side of musculoskeletal contraction and movement, including proprioception, tactile feedback and other task-relevant sensations linked to action. It marks proposed scope, without establishing a new external scientific standard. That scope supplies neither acquired coverage nor body-wide simultaneous observations, and the four groups are not an exhaustive taxonomy.

Known tactile stimuli, proprioceptive feedback, contact, force and task context contribute different evidence. Stimulus delivery, reported sensation, recorded sensory activity, calibrated force and tracked configuration cannot be substituted without a tested observation relationship. EMG records muscle electrical activity, not direct peripheral nerve or sensory activity; kinematics alone does not measure proprioceptive experience. Each measured, reported, estimated or simulated quantity keeps its provenance.

The simulation group supplies potential training examples as well as constraints: trajectories, state transitions, muscle--tendon or contact-force outputs and declared movement priors. Record the simulator, body-model version, parameters, solver, control representation and underlying sources. Joint-target control must not be called muscle-driven control. Retargeted recordings retain original lineage and transformations; generated states or forces remain simulated. The generator requires calibration against appropriate measured support with uncertainty and training-only role boundaries.

Test its contribution by comparing a fixed-real-data reference with a simulation-augmented arm under matched permitted information, compute and original criteria. Record real-example exposure, simulated exposure and updates in each arm; a fixed real corpus alone does not match training exposure. Count simulated amount and diversity separately from real acquisitions. Held-out real observations must check the result without entering generator calibration, tuning, reward choice or model selection. Wrong dynamics, contacts or source correspondence can make the synthetic contribution fail. Generated data is never independent biological ground truth for the mechanism that produced it. Generator methods, implementation and use conditions remain to be checked; no frontier implementation or dataset is selected here.

Five component roles describe another axis: brain signals, muscle signals, measured movement, a virtual body and a learning model. A virtual body can supply training data and physical constraints, while a learning model can draw from several source groups. Neither count excludes another useful, properly characterized observation. These resources give the later teaching study more hypotheses to test, rather than supplying the missing retained-learning result.

\subsection{Breadth needs a controlled comparison}

Amount and diversity must be varied separately. Repeated windows from one acquisition increase presentations without adding independent subjects, days or tasks. A proposed amount comparison would increase support from a stated source distribution. A diversity comparison would retain a declared amount while changing named aspects of that distribution, with a crossed comparison testing their interaction. Organism, subject, session, movement, body region, modality, apparatus, environment and measurement quality require separate descriptions. Dataset count alone cannot measure them.

The architecture, permitted input information, targets, preprocessing and selection rules would remain fixed. Source rosters, repetitions, updates and measured computing work must accompany the results; equal update counts alone do not establish matched compute. If diversity also changes recording quality or task difficulty, the comparison cannot attribute its outcome to diversity alone. All exact budgets, controls, seeds, uncertainty methods and decisions remain to be registered before a new experiment. This proposal changes neither the adopted continuous WAY-EEG-GAL development design nor any original protected partition.

A candidate amount--error curve, for a fixed diversity condition and evaluation domain, is
\[
E(N)=E_{\infty}+A(N/N_0)^{-\alpha}.
\]
$N>0$ and its fixed reference $N_0>0$ use the same declared training-amount units. $E$ is a chosen nonnegative error; $E_{\infty}$ and $A$ use its units, while $\alpha$ is dimensionless. This decreasing-error candidate requires $E_{\infty}\geq0$, $A\geq0$ and $\alpha>0$: $A>0$ gives strict decrease, while $A=0$ gives constant error. Positive $\alpha$ alone is insufficient. Predictions must also respect any upper bound of the metric throughout the declared fitting and prediction range. The floor is a prospective curve parameter, not an established irreducible biological limit. No exponent is estimated here and no universal law follows. Constant, saturating, increasing, non-monotonic and domain-specific alternatives need prospective comparison and predictive checks at withheld scaling levels. Observed negative transfer remains a result even when this constrained candidate cannot fit it. Related windows cannot supply independent uncertainty replicates. Neural-target, EMG-target, pixel and physical-motion errors retain their own units and denominators.

Generalization would be evaluated separately across excluded subjects, sessions, tasks, datasets and, where mechanically qualified, species. The target-only/no-transfer reference, applicable incorrect correspondence, simpler predictors and original null, shuffle, cue/confound and causal-time controls remain. Negative transfer, worse domains and saturation stay in the record. A better pooled score cannot discharge an original per-output no-harm criterion or erase known pretraining exposure.

\subsection{Build a composite without inventing a shared observation}

FALCON M1-A and AJILE12 illustrate complementary source candidates. The fixed FALCON M1-A release describes simultaneous M1 activity and processed intramuscular EMG from one rhesus macaque; released synchronized motion capture is not established. Its acausally processed EMG targets limit causal real-time interpretation. The small minival subset is drawn from calibration data for submission-format checks and is not independent validation. AJILE12 describes intracranial recordings and two-dimensional, video-estimated upper-body pose from twelve clinical participants. Paired EMG and instrumented three-dimensional mocap are not established there. Its pose-label construction is participant-adapted; particular held-out-day estimator exposure and numerical alignment uncertainty remain unqualified. \href{https://doi.org/10.48324/dandi.000941/0.241029.1405}{FALCON M1-A release}; \href{https://papers.nips.cc/paper_files/paper/2024/file/8c2e6bb15be1894b8fb4e0f9bcad1739-Paper-Datasets_and_Benchmarks_Track.pdf}{FALCON methods}; \href{https://doi.org/10.1038/s41597-022-01280-y}{Peterson et al., 2022}; \href{https://doi.org/10.48324/dandi.000055/0.220127.0436}{AJILE12 release}.

Qualified source-specific pretraining could exploit each relationship. A transfer model would then have to demonstrate a correspondence on actual target measurements. The separate recordings never become one simultaneous neural--muscle--movement event. The release descriptions do not establish downloaded recordings or a trained model. Their exact use conditions, versioned files, integrity, original acquisition relationships, timing, units, earlier exposure and independent partitions still need checking before training. Retargeting between homologous or mechanically analogous structures also requires declared geometry, frames, scale, load, contact, timing and uncertainty. Similar movement cannot establish identical neural coding.

\subsection{Model the inside and the outside of one action}

When a hand closes on an object, external displacement and contact put demands on internal mechanics. The proposed world model would represent specified musculoskeletal states alongside the supported body, objects and environment. NeuroTranslate would connect named observations and targets; MoCap ReComp would provide its separately qualified source-preserving movement composite; MyoSim Replay would simulate declared mechanics. Their proposed combination is the Large Movement Model. Neither that name nor a neural composite implies a measured whole nervous system, equivalent cells across organisms, complete internal organs or unique access to private experience.

The model must distinguish measured neural activity, muscle electrical activity and tracked motion from estimated activation, internal forces, latent neural state and simulated motion. EMG is not direct peripheral nerve recording or automatically calibrated force. A proposed transition model could forecast a distribution over joint configuration, generalized velocity, muscle-model activation and environmental state from permitted past observations, context and declared controls. Separate observation models would connect each sensor to the state it can constrain. An unavailable observation stays unavailable; a conditional reverse prediction cannot identify a unique hidden brain.

Physical tests would compare aligned predicted movement with independent measured motion, and any internal-force claim with appropriate independent force evidence. Joint limits, contacts, loads, energy accounting, friction, gravity and numerical solver behaviour need their own checks. OpenSim supplies established methods for declared musculoskeletal models; it cannot confer biological validity on this proposed assembly. \href{https://journals.plos.org/ploscompbiol/article?id=10.1371/journal.pcbi.1006223}{Seth et al., 2018}. A rendering can show where a model says a muscle pulls. Only the applicable observations can test whether that pull explains the claimed result.

% R4 conceptual-article insertion candidate; no manuscript/source promotion.
% Inherit the selected master's Noto Sans/Noto Sans Mono/Fira Math setup.
\subsection{One muscle pulse, two simulated trajectories}

A moving virtual hand has answered a mechanical question before it has answered a teaching question. In one numerical diagnostic, we compared two trajectories of the installed \texttt{myohand\_r} human right-hand mechanism on a fixed passive torso. A muscle excitation was added to one branch. The other supplied its matched control. This asks whether the specified software input changes the computed state. It provides no neural decoding, measured biological response or retained-learning outcome. The mechanism used \texttt{myo\_sim} 0.2.3, commit \texttt{93b0ca8f4ec90c9899ee7f05fee561e9911da91b}, MuJoCo 3.14.0 and NumPy 2.5.3.\cite{myohandMyoSimPinned2026,myohandMuJoCo3140,mujoco2012,myohandNumericalRecord20261105}

\subsubsection{Mechanism and matched comparison}

The compiled model has 23 scalar hinge degrees of freedom, 39 muscle--tendon actuators, 39 activation states and zero sensors. Both branches began with matching full recorded initial states, including positions, velocities and activation states. The control branch received no additional excitation pulse; the pulse branch received dimensionless excitation 0.2 at \texttt{FDS2\_r}, with all other controls zero. Excitation is a muscle-model input, distinct from a joint-position request, measured EMG or neural activity. Gravity and passive dynamics remain present, so the unpulsed trajectory need not remain stationary.

The native Euler integrator advanced each branch through 100 explicit steps of $0.002\,\mathrm{s}$, giving nominal $T=0.200\,\mathrm{s}$. Pulse steps 10--59 inclusive correspond to the pre-step interval $0.020\leq t<0.120\,\mathrm{s}$. Each branch received explicit initial and final forward evaluations to synchronize state-dependent quantities. Thus the record contains 200 explicit integration steps and four explicit forward evaluations. The executed source path also implies two compilation calls and one construction-time internal forward evaluation. Those counts do not enumerate every internal native solver operation.

The construction path changed physical parameters even with optional \texttt{inertia\_floor=None}: each diagonal inertia component of four fixed scaffold bodies changed from $1$ to $10^{-4}\,\mathrm{kg\,m^2}$. This native transformation is part of the tested mechanism, rather than evidence that all parameters remained untouched. No additional physical retuning or \texttt{implicitfast} fallback was used. Position, orientation, scale, contact behaviour and biological fidelity have not been independently calibrated by this diagnostic.

\subsubsection{Recorded response}

For pulse branch $P$, control $C$ and the synchronized final state, the reported quantities are
\[
\Delta q_{\infty}=\max_{1\leq j\leq23}|q_{P,j}(T)-q_{C,j}(T)|,
\qquad
\Delta a=a_{P,\mathrm{FDS2}}(T)-a_{C,\mathrm{FDS2}}(T),
\]
\[
\Delta F=F_{P,\mathrm{FDS2}}(T)-F_{C,\mathrm{FDS2}}(T).
\]
The first quantity is the maximum absolute separation across the 23 hinge coordinates at the final time, in radians. The other quantities are signed pulse-minus-control contrasts in the selected dimensionless activation and tendon-actuator force in newtons. Force follows the model's convention and uses synchronized endpoint forward evaluations, rather than derived values left over from the last integration step.

% Shared by all four proposed R3 inserts; no added packages or font defaults.
\begin{table}[htbp]
\centering
\caption{Final contrasts for one simulated right-hand mechanism. All values compare the pulse branch with its matched control at nominal $T=0.200\,\mathrm{s}$. The angle value is the largest absolute separation across 23 hinges at that endpoint; activation and force retain the pulse-minus-control sign. Activation and force concern the selected \texttt{FDS2\_r} actuator. Force uses the synchronized final forward evaluations and the model's sign convention. These are simulated contrasts, not errors against biological measurements. The displayed digits preserve the recorded numerical values, without establishing biological accuracy.\cite{myohandNumericalRecord20261105}}
\label{tab:myohand-numerical-response}
\begin{tabular}{p{0.53\linewidth}p{0.26\linewidth}p{0.10\linewidth}}
\hline
Quantity & Contrast & Unit \\
\hline
Largest absolute final joint-coordinate separation, $\Delta q_{\infty}$ & $0.5027347840656727$ & rad \\
Selected final activation difference, $\Delta a$ & $0.05205729904534248$ & none \\
Selected final tendon-actuator force difference, $\Delta F$ & $-6.407722136774943$ & N \\
\hline
\end{tabular}
\end{table}


Independent auditing of the saved evidence completed 132 standard-library checks without a failed check or model rerun. All required recorded numerical channels were finite and recorded native warning counts were zero. Four stored endpoint mass-matrix spectra were positive. Raw mass matrices and residual operands were not independently exported; their stored diagnostics were checked without independently resolving those matrices. No contact was observed. At the final control endpoint, the recorded soft-limit excess was $0.0004121829780273552\,\mathrm{rad}$, so exact hard-bound enforcement is not established. Neither conservation, numerical convergence nor a physiological response threshold follows.

These contrasts establish a response to the added excitation within this simulated configuration. They are not errors against measured human movement, EMG or force. Two dependent trajectories with 200 time steps supply neither 200 independent experiments nor biological replication. No statistical $p$-value is inferred, and numerical precision in the record does not establish biological accuracy or repeatability across cases.

\subsubsection{From a response to a teaching experiment}

The diagnostic establishes the recorded input-to-state response, not a performance benchmark. Its complete computing-resource record and unresolved memory/timing attribution are retained in Appendix~\ref{app:reproducibility-availability}.

The hand's zero sensors leave Myoception---the programme's proposed sensory side of muscle-driven action---without a measurement in this test. Simulated activation, force and motion do not supply afferent recordings, biological feedback, a learned controller, a trained Large Movement Model or generalization. The four proposed source groups remain neural observations, myography, measured motion and physics-generated examples. Mechanics, controls, calibration, uncertainty, source rights and data lineage require qualification before this generator could enter training, and its contribution would need independent held-out real tests. The amount/diversity hypothesis remains untested.

For Brain2Body, even a later qualified controller would leave the teaching question: what does a trainee retain when trainer commands end? That comparison still needs the original matched information, practice and assistance controls, sham/delayed/unrelated guidance, independent outcomes and a constant compensatory movement bridge where required. This diagnostic contains no trainee, teaching, withdrawal or physical intervention. It selects no training generator and supplies no biological ground truth. LINK's eighteen fits and 720 saved epochs retain both failed admission gates and all five failed benefit claims; protected tests and quarantined sources remain closed.

\subsection{Carry the test through teaching and withdrawal}

A better movement prior might improve assistance, but the teaching hypothesis requires another comparison. A prospective trainer--trainee study must identify the model's inputs, expert information, feedback, commands, practice and assistance burden. Match those resources against the appropriate simpler controller and original sham, delayed and unrelated-guidance controls. Otherwise an apparent neural benefit could come from supplying the target action, more practice or a less demanding task.

Retention is assessed after trainer commands stop, with any necessary compensatory movement bridge held constant. Immediate performance, delayed retention and unfamiliar task variations remain separate outcomes. A gain that depends on continued guidance may support assistance. A controlled gain that persists would need its own domain and uncertainty; it would still establish neither tissue repair nor arbitrary skill transfer. Simulation, offline replay, measured biological feedback and physical intervention retain distinct evidence and authority.

No scaling result, trained body-wide world model or biological teaching result is reported here. LINK retains eighteen fits and 720 saved epochs, with both model-admission gates and all five benefit claims failed. A broader programme gives those failures somewhere useful to lead: a new, source-qualified comparison that can show exactly which prediction improves, what supplies the improvement and whether a trainee retains anything from it.
